% d02-xcolor: xcolor colours in text, boxes and rules (RGB, CMYK, gray, named, HTML, mixes), colorbox, fcolorbox, coloured text across a page break
\documentclass{article}
\usepackage{fontspec}
\usepackage[dvipsnames]{xcolor}
\definecolor{mine}{HTML}{3A7F2B}
\colorlet{soft}{blue!30!white}
\begin{document}
Black \textcolor{red}{red \textcolor{blue}{blue} red} black.
{\color[cmyk]{0.1,0.8,0.2,0.05} CMYK text} and {\color[gray]{0.5} gray text}
and {\color[RGB]{200,100,0} RGB 0-255} and \textcolor{mine}{HTML} and \textcolor{soft}{colorlet}.
\colorbox{yellow}{boxed} \fcolorbox{OliveGreen}{Peach}{framed} \colorbox{soft}{\textcolor{white}{white on soft}}
{\color{RoyalBlue!60!black}\rule{3cm}{2pt}} \textcolor{green!40!black}{mix}
\textcolor{red!50}{tint} \textcolor{-red}{complement} {\color{gray}\rule{1cm}{1cm}}
\rule{4cm}{0.4pt} {\color{red}\rule[-1pt]{0.5pt}{8pt}}
\begin{center}\color{violet}Centered violet \textbf{bold}\end{center}
\vspace*{\fill}
{\color{Bittersweet}This paragraph is set in one colour and runs over the page break so the
colour has to be restored at the top of the next page. It needs enough words to reach the bottom of
the page and continue on the next one, so here are some more words to make it long enough to break.
And a few more, and more again, to be sure the paragraph crosses the boundary between the two pages
and the colour stack is carried over. \textcolor{blue}{Nested blue text inside, also crossing perhaps,
with yet more words so that it spans lines.} Back to the outer colour for the rest of this paragraph,
which ends here after a few more words.\par}
Normal black text after the coloured paragraph.
\end{document}
